\section{Machines and observable sections}\label{sec:model}
We use row vectors for state distributions. Put
$\Delta_k=\{e\in\R^k:e\ge0,\ e\one=1\}$,
$S=[-1/10,1/10]^2$, and $D=[-1,1]^2$.

\begin{definition}[Clocked realization]\label{def:machine}
A realization of \eqref{eq:array} consists of an initialization
$E_x\in\Delta_{K_0}$, row-stochastic matrices
$T_{t,c}\in\R_+^{K_t\times K_{t+1}}$ for $0\le t<N$, $c\in\{0,1\}$,
and decoder means $d_j\in[-1,1]^{K_N}$, satisfying
\begin{equation}\label{eq:machine}
 E_xT_{0,c_1}\cdots T_{N-1,c_N}d_j
                  =e_j^TR_\alpha^{\sum c_t}x/10
\end{equation}
for every input word and query. The decoder emits the binary law with
this conditional mean. The available-label profile is $(K_t)_{t=0}^N$;
its peak is $\max_tK_t$, and its minimum is $W_N(\alpha)$.
A final retained binary output can be implemented in two further output
labels at its own cut and does not change this peak, which is at least three.
\end{definition}

Only the current register persists. The seed, old commands, random coins,
and protocol history are not supplied again. The seed and terminal query
are atomic inputs; commands are external, rather than selected by an
experimental policy. A row may depend on the external epoch and current
command, but not on a future symbol. Every row, including an unreachable
one, is normalized. Available-label counting permits padding and makes
feasible profiles upward closed. The number of program bits, exact row
arithmetic, and sampling cost are not part of this width.
For comparison, a separate minimum at cut $t$ imposes no bound at any
other cut. It is not the peak of a simultaneously minimal machine.

\begin{definition}[Vector-state realization]\label{def:vector}
A clocked realization is vector-state if there are vectors
$z_{t,s}\in D$ satisfying
\begin{equation}\label{eq:statewise}
 \sum_s E_x(s)z_{0,s}=x/10,\qquad
 \sum_{s'}T_{t,c}(s,s')z_{t+1,s'}=R_\alpha^c z_{t,s},
 \qquad d_j(s)=e_j^Tz_{N,s}.
\end{equation}
The middle equality is required at every available state, not only after
averaging over reachable distributions. The least peak in this subclass
is $V_N(\alpha)$.
\end{definition}

Such a machine is equivalent to a sequence of polygons
$P_t=\conv\{z_{t,s}\}$ with
\begin{equation}\label{eq:chain}
 S\subset P_0,\qquad P_t\subset D,\qquad
                 \conv(P_t\cup R_\alpha P_t)\subset P_{t+1}.
\end{equation}
The forward implication follows from \eqref{eq:statewise}; conversely,
choose the vertices as states and the corresponding barycentric rows.
Thus the minimum vertex count of a compatible chain is exactly the
vector-state optimization. No such identification is imposed on $W_N$.

\begin{proposition}[Separate minimum]\label{prop:separate}
Every internal cut of \eqref{eq:array} has normalized residual rank
three and separate positive minimum three, for every $N$ and rotation.
\end{proposition}
\begin{proof}
The four prefixes with all commands idle have predictive vectors $x/10$
and span an affine plane. The two idle-suffix coordinate queries separate
this plane. All residual arrays are affine functions of a two-dimensional
predictive vector, so their normalized rank is exactly three.
The vectors have Euclidean norm $\sqrt2/10$. The triangle with vertices
\[
             (1/2,0),\quad(-1/2,1/2),\quad(-1/2,-1/2)
\]
contains the disk of this radius: its three supporting lines have
distances $1/2,1/(2\sqrt5),1/(2\sqrt5)$ from zero, all greater than
$\sqrt2/10$. Its vertices have norm at most $1/\sqrt2$.
Consequently every vertex defines a legal entire future continuation
under all rotations and both queries. Barycentric initialization at the
selected cut gives a three-generator factorization of all its residuals.
To embed it into a complete machine, store the raw prefix until the
selected cut, sample that three-valued label, and then retain the label
and the remaining finite command count. Other cuts are unrestricted.
Rank supplies the matching lower bound. This is the inherited separate
minimum calculation of \cite{GTF32}, included to fix its quantifier.
\end{proof}

\begin{lemma}[Reachable-section lift]\label{lem:reachable}
Every realization determines affine sections
$Q_t=H_t\cap\Delta_{K_t}$ and linear maps $\pi_t$ with
$P_t=\pi_t(Q_t)$ satisfying \eqref{eq:chain} and
\begin{equation}\label{eq:face-count}
                         |\ext P_t|\le2^{K_t}.
\end{equation}
For vector-state realizations one may instead use $|\ext P_t|\le K_t$.
\end{lemma}
\begin{proof}
Let $E_h$ denote the state distribution following an actual prefix, and
set $H_t=\aff\{E_h\}$. Define $\pi_t(e)$ by the two terminal means
obtained from $e$ after idling all remaining commands. It maps the whole
simplex into $D$. The identities
\[
 E_hT_{t,c}=E_{hc},\qquad
 \pi_{t+1}(E_hT_{t,c})=R_\alpha^c\pi_t(E_h)
\]
extend affinely across $H_t$. Stochasticity then gives
$Q_tT_{t,c}\subset Q_{t+1}$ and the required observable inclusions.
The idle command implies $P_t\subset P_{t+1}$, hence every $P_t$
contains $S$.
In its affine hull $Q_t$ has at most $K_t$ coordinate inequalities.
Distinct vertices have distinct active zero-coordinate sets: otherwise
their difference is a two-sided feasible direction at either vertex.
There are at most $2^{K_t}$ such sets, and projection cannot increase
the number of vertices. This is the standard face-count bound for
extensions \cite{FRT}, not a special new combinatorial principle.
For the vector-state class, use the polygon in \eqref{eq:chain}.
\end{proof}

The restriction to $H_t$ is essential. A hidden basis state need not have
the observable intertwining identity. The proof therefore retains the
section before projection, as in accessible-space and observable-quotient
constructions \cite{BF04,MW}. It is this lift that separates the two
bounds in Lemma~\ref{lem:reachable}.
